\documentclass[a4paper,10pt]{article}

\usepackage[T1]{fontenc}
\usepackage[utf8]{inputenc}
\usepackage[margin=0.65in]{geometry}
\usepackage{titlesec}
\usepackage{graphicx}
\usepackage{xcolor}
\usepackage{enumitem}
\usepackage[hidelinks]{hyperref}
\usepackage{tabularx}
\usepackage{fancyhdr}
\usepackage{array}

\pagestyle{fancy}
\fancyhf{}
\renewcommand{\headrulewidth}{0pt}
\renewcommand{\footrulewidth}{0pt}

\setlength{\parindent}{0pt}
\setlength{\parskip}{3pt}
\setlength{\tabcolsep}{0pt}
\raggedbottom
\raggedright

\definecolor{accent}{HTML}{4F46E5}
\definecolor{accent2}{HTML}{14B8A6}
\definecolor{textmain}{HTML}{111827}
\definecolor{textsoft}{HTML}{4B5563}

\titleformat{\section}{\vspace{-2pt}\scshape\large\bfseries\color{accent}}{}{0em}{}[\color{accent2}\titlerule\vspace{-4pt}]

\pdfgentounicode=1

\graphicspath{{old_hag/assets/}{../old_hag/assets/}{public/}{../public/}{src/assests/logos/}{../src/assests/logos/}{./}}

\newcommand{\resumeItem}[1]{\item\small{#1\vspace{-2pt}}}
\newcommand{\resumeSubheading}[4]{
  \vspace{-2pt}\item
    \begin{tabular*}{0.98\textwidth}[t]{l@{\extracolsep{\fill}}r}
      \textbf{#1} & \textbf{#2} \\
      \textit{\small#3} & \textit{\small #4} \\
    \end{tabular*}\vspace{-7pt}
}
\newcommand{\resumeProjectHeading}[2]{
  \item
    \begin{tabular*}{0.98\textwidth}{l@{\extracolsep{\fill}}r}
      \small\textbf{#1} & \small #2 \\
    \end{tabular*}\vspace{-6pt}
}
\newcommand{\resumeSubHeadingListStart}{\begin{itemize}[leftmargin=0.15in,label={}]}
\newcommand{\resumeSubHeadingListEnd}{\end{itemize}}
\newcommand{\resumeItemListStart}{\begin{itemize}[leftmargin=*,itemsep=2pt,topsep=2pt]}
\newcommand{\resumeItemListEnd}{\end{itemize}\vspace{-4pt}}

\begin{document}

\begin{minipage}[c]{0.18\textwidth}
  \includegraphics[width=\linewidth,height=3.6cm,keepaspectratio]{Photo.png}
\end{minipage}
\hfill
\begin{minipage}[c]{0.78\textwidth}
  \begin{center}
    {\LARGE \textbf{\textcolor{textmain}{Walid Kessoum}}}\\[3pt]
    {\large \textcolor{accent}{Ingénieur Réseaux \& Télécommunications \;\textbullet\; Performance \& Optimisation}}\\[8pt]
    {\small Paris, Île-de-France $|$ +33\,7\,46\,50\,54\,37 $|$ \href{mailto:mw.kessoum@icloud.com}{mw.kessoum@icloud.com} $|$ \href{https://linkedin.com/in/walid-kessoum}{LinkedIn} $|$ \href{https://propanone.github.io/}{Portfolio}}\\[8pt]
    {\small Ingénieur Télécoms \& Réseaux, double diplôme Algérie/France, orienté ingénierie opérateur, performance radio, transport et fiabilité réseau. Expérience concrète sur les chaînes RAN/Core, les KPI LTE, l'automatisation data et les architectures multi-sites. Lauréat mondial Huawei ICT.}
  \end{center}
\end{minipage}

\vspace{4pt}

\section{Compétences Réseaux \& Télécoms}
\begin{itemize}[leftmargin=0.16in,label={}]
  \small{\item{
    \textbf{Réseaux mobiles \& opérateur}{: LTE/5G, RAN \& Core, KPI radio, TCP/IP, VoIP, FTTH, SDH, DWDM, OFDM, OSPF multi-area, NAT, DHCP, Cisco Packet Tracer} \\
    \textbf{Performance \& exploitation}{: Data Warehouse, ETL, Apache Airflow, Python, SQL/PostgreSQL, LSTM, PowerBI, monitoring temps réel, diagnostic réseau} \\
    \textbf{Écosystème \& outils}{: Wireshark, Linux, Docker, Git, Huawei Cloud, Huawei ICT, Cisco, OTDR, SDIP / certification data-driven} \\
    \textbf{Développement}{: Django, Flask, C/C++, MATLAB, HTML/CSS} \\
    \textbf{Langues}{: Français (courant) \; $\cdot$ \; Anglais (courant) \; $\cdot$ \; Arabe (natif)
  }}
\end{itemize}

\section{Formation}
\resumeSubHeadingListStart
  \resumeSubheading
    {Université Paris-Saclay \& Institut Polytechnique de Paris}{Gif-sur-Yvette, France}
    {Master 2 Multimedia Networking}{Sept. 2026 -- Sept. 2027}
  \resumeSubheading
    {Sup Galilée -- Université Sorbonne Paris Nord}{Paris, France}
    {Cycle Ingénieur (2\textsuperscript{ème} année), Spécialité Réseaux et Télécommunications}{Sept. 2025 -- Juin 2026}
  \resumeSubheading
    {ENSTA -- École Nationale Supérieure de Technologie Avancée}{Alger, Algérie}
    {Ingénieur d'État \& Master en Systèmes de Télécommunications et Réseaux}{Oct. 2020 -- Juin 2025}
    \resumeItemListStart
      \resumeItem{Projet de Fin d'Études : \textit{A Data-Driven Framework for AI-Powered LTE Network Performance Optimization} -- mention très bien}
    \resumeItemListEnd
\resumeSubHeadingListEnd

\section{Expérience Professionnelle}
\resumeSubHeadingListStart
  \resumeSubheading
    {Djezzy}{Oct. 2024 -- Mai 2025}
    {Ingénieur Data \& Performance Réseau (Alternance)}{Alger, Algérie}
    \resumeItemListStart
      \resumeItem{Conçu un Data Warehouse PostgreSQL centralisant les données LTE (RAN \& Core) avec pipelines ETL automatisés via Python et Apache Airflow}
      \resumeItem{Développé des modèles LSTM pour la prévision de KPIs critiques et l'aide à la détection d'anomalies}
      \resumeItem{Créé une plateforme de supervision fullstack (Django/Flask + PowerBI) pour le suivi opérationnel réseau}
    \resumeItemListEnd

  \resumeSubheading
    {Huawei ICT Competition}{Mai 2025}
    {Lauréat Mondial (3\textsuperscript{ème} Prix) -- Innovation Track}{Shenzhen, Chine}
    \resumeItemListStart
      \resumeItem{Développé une solution SaaS d'IA pour l'assurance avec XGBoost sur Huawei Cloud ModelArts}
      \resumeItem{Classé 3\textsuperscript{ème} parmi 150\,000+ participants mondiaux après un pitch technique en anglais devant jury international}
    \resumeItemListEnd

  \resumeSubheading
    {SONATRACH}{Fév. 2024 -- Avr. 2024}
    {Analyste Réseau (Stage)}{Alger, Algérie}
    \resumeItemListStart
      \resumeItem{Analyse des architectures SCADA et VoIP et étude des contraintes de disponibilité et de performance}
      \resumeItem{Automatisation de tâches réseau avec scripts Python pour accélérer le diagnostic}
    \resumeItemListEnd

  \resumeSubheading
    {Algérie Télécom}{Mars 2023 -- Avr. 2023}
    {Assistant Ingénieur Réseaux (Stage)}{Alger, Algérie}
    \resumeItemListStart
      \resumeItem{Déploiement d'infrastructures FTTH avec installation fibre optique et diagnostics OTDR}
      \resumeItem{Exposition aux technologies SDH, DWDM, OFDM et gestion du trafic au niveau transport}
    \resumeItemListEnd
\resumeSubHeadingListEnd

\section{Projets Sélectionnés}
\resumeSubHeadingListStart
  \resumeProjectHeading
    {\textbf{SWITCHERS -- Architecture Réseau d'Entreprise Résiliente (Cisco Packet Tracer)}}{Sup Galilée}
    \resumeItemListStart
      \resumeItem{Conception d'une architecture multi-sites avec OSPF multi-area, DHCP, HTTP/Mail, NAT et accès Internet sécurisé}
      \resumeItem{Simulation de panne de lien avec reconvergence automatique et validation de la continuité de service}
    \resumeItemListEnd

  \resumeProjectHeading
    {\textbf{Data Warehouse LTE -- Plateforme centralisée de performance réseau}}{Projet de recherche / mémoire}
    \resumeItemListStart
      \resumeItem{Plateforme d'agrégation et de monitoring des KPIs LTE avec pipelines data automatisés et modèles de prévision}
      \resumeItem{Approche orientée exploitation opérateur, industrialisation ETL et visualisation PowerBI}
    \resumeItemListEnd

  \resumeProjectHeading
    {\textbf{SCADA/VoIP Network Analysis}}{Projet réseau industriel}
    \resumeItemListStart
      \resumeItem{Étude de fiabilité, latence et contraintes de diagnostic sur un environnement réseau industriel}
    \resumeItemListEnd
\resumeSubHeadingListEnd

\section{Certifications \& Distinctions}
\resumeSubHeadingListStart
  \resumeProjectHeading
    {\textbf{Huawei ICT Competition -- Global Winner (3\textsuperscript{ème} Prix)}}{2025}
    \resumeItemListStart
      \resumeItem{Reconnaissance mondiale pour excellence technique en innovation réseau et cloud}
    \resumeItemListEnd

  \resumeProjectHeading
    {\textbf{Google Cybersecurity, Google IT Automation with Python, Google IT Support}}{2022 -- 2024}
    \resumeItemListStart
      \resumeItem{Certifications complémentaires renforçant le socle cybersécurité, scripting et support systèmes}
    \resumeItemListEnd
\resumeSubHeadingListEnd

\end{document}